\section{Autoduality of the Planck internal unit under orientation inversion}

The CT108 selector constructs

\begin{equation}
\theta_{108}=\left(\frac{54}{\pi}\right)^2,
\qquad
\ell_P=\frac{54}{\pi}\ell_0,
\qquad
t_P=\frac{54}{\pi}t_0,
\label{eq:planck-length-time}
\end{equation}

and

\begin{equation}
E_P=\frac{\hbar_{\rm int}}{t_P}
=\frac{\pi\hbar_{\rm int}}{54t_0}.
\label{eq:planck-energy}
\end{equation}

Comparing \cref{eq:boltzmann-clock,eq:planck-energy} gives

\begin{equation}
\boxed{E_P=k_B\Theta_{\rm clk}\log3.}
\label{eq:planck-trit}
\end{equation}

The relation does not depend on a coincidence between units: both sides
follow from the same chronological structure \(108\).

Once the gravitational constant has been derived internally in the
\cref{ch:gravity}, the Planck family satisfies

\begin{align}
\ell_P^2&=\frac{G\hbar}{c^3}=\theta_{108}\ell_0^2,
\label{eq:planck-area}\\
F_P&=\frac{c^4}{G},
\qquad
P_P=\frac{c^5}{G},
\label{eq:planck-force-power}\\
\rho_P&=\frac{\hbar}{\theta_{108}^2c^5t_0^4}.
\label{eq:planck-density}
\end{align}

Although \(G\) and \(\hbar\) transform reciprocally between the sections
before and after the return, their product in \(\ell_P^2\) preserves the
geometry. This invariance determines area as the interface between the
geometric, informational and entanglement realizations.
